\SourceLeaf{022}
gradus est potestas binaria, nihil doceat, insuperque omnibus obiectionibus praecc.
(praeter quartam) aeque obnoxia sit ut demonstratio prima generalis: haud necesse
est illam hic fusius explicare.

\Article{9}
In eadem commentatione ill. E. theorema nostrum adhuc alia via confirmare
annixus est p 263, cuius summa continetur in his: Proposita aequatione
$x^n+Ax^{n-1}+Bx^{n-2}\ \text{etc.}=0$, hucusque quidem expressio analytica, quae ipsius
radices exprimat, inveniri non potuit, si exponens $n>4$; attamen certum esse
videtur (uti asserit E.), illam nihil aliud continere posse, quam operationes arithmeticas
et extractiones radicum eo magis complicatas, quo maior sit $n$. Si hoc
conceditur, E. optime ostendit, quantumvis inter se complicata sint signa radicalia,
tamen formulae valorem semper per formam $M+N\BareRoot{-1}$ repraesentabilem
fore, ita ut $M,N$ sint quantitates reales.

Contra hoc ratiocinium obiici potest, post tot tantorum geometrarum labores
perexiguam spem superesse, ad resolutionem generalem aequationum algebraicarum
umquam perveniendi, ita ut magis magisque verisimile fiat, talem resolutionem
omnino esse impossibilem et contradictoriam. Hoc eo minus paradoxum
videri debet, \textit{quum id, quod vulgo resolutio aequationis dicitur, proprie nihil aliud sit
quam ipsius reductio ad aequationes puras.} Nam aequationum purarum solutio hinc
non docetur sed supponitur, et si radicem aequationis $x^m=H$ per $\BareRoot[m]{H}$ exprimis,
illam neutiquam solvisti, neque plus fecisti, quam si ad denotandam radicem
aequationis $x^n+Ax^{n-1}+\text{etc.}=0$ signum aliquod excogitares, radicemque huic
aequalem poneres. Verum est, aequationes puras propter facilitatem ipsarum radices
per approximationem inveniendi, et propter nexum elegantem, quem omnes
radices inter se habent, prae omnibus reliquis multum praestare, adeoque neutiquam
vituperandum esse, quod analystae harum radices per signum peculiare denotaverunt:
attamen ex eo, quod hoc signum perinde ut signa arithmetica additionis,
subtractionis, multiplicationis, divisionis et evectionis ad dignitatem sub
nomine \textit{expressionum analyticarum} complexi sunt, minime sequitur cuiusvis aequationis
radicem per illas exhiberi posse. Seu, missis verbis, sine ratione sufficienti
supponitur, cuiusvis aequationis solutionem ad solutionem aequationum purarum
reduci posse. Forsan non ita difficile foret, impossibilitatem iam pro quinto gradu
omni rigore demonstrare, de qua re alio loco disquisitiones meas fusius proponam.
\SourceLeaf{023}
Hic sufficit, resolubilitatem generalem aequationum, in illo sensu acceptam. adhuc
valde dubiam esse, adeoque demonstrationem, cuius tota vis ab illa suppositione
pendet, in praesenti rei statu nihil ponderis habere.

\Article{10}
Postea etiam clar. \textsc{de Foncenex}, quum in demonstratione prima \textsc{Euleri} defectum
animadvertisset (supra art. 8 obiect. 4), quem tollere non poterat, adhuc
aliam viam tentavit et in comment. laudata p. 120 in medium protulit\HistoricalNote{7}.
Quae consistit in sequentibus.

Proposita sit aequatio $Z=0$, designante $Z$ functionem $m^{\mathrm{ti}}$ gradus incognitae
$z$. Si $m$ est numerus impar, iam constat, aequationem hanc habere
radicem realem; si vero $m$ est par, clar. F. sequenti modo probare conatur, aequationem
ad minimum unam radicem formae $p+q\BareRoot{-1}$ habere. Sit $m=2^n i$,
designante $i$ numerum imparem, supponaturque $zz+uz+M$ esse divisor
functionis $Z$. Tunc singuli valores ipsius $u$ erunt summae binarum radicum aequationis
$Z=0$ (mutato signo), quamobrem $u$ habebit $\frac{m\mathbin{.}m-1}{1\mathbin{.}2}=m'$ valores,
et si $u$ per aequationem $U=0$ determinari supponitur (designante $U$ functionem
integram ipsius $u$ et coëfficientium cognitorum in $Z$), haec erit gradus $m'^{\mathrm{ti}}$.
Facile vero perspicitur, $m'$ fore numerum formae $2^{n-1}i'$, designante $i'$ numerum
imparem. Iam nisi $m'$ est impar, supponatur iterum, $uu+u'u+M'$ esse divisorem
ipsius $U$, patebitque per similia ratiocinia, $u'$ determinari per aequationem $U'=0$,
ubi $U'$ sit functio $\frac{m'\mathbin{.}m'-1}{1\mathbin{.}2}^{\mathrm{ti}}$ gradus ipsius $u'$. Posito vero $\frac{m'\mathbin{.}m'-1}{1\mathbin{.}2}=m''$, erit $m''$
numerus formae $2^{n-2}i''$, designante $i''$ numerum imparem. Iam nisi $m''$ est impar,
statuatur $u'u'+u''u'+M''$ esse divisorem functionis $U'$, determinabiturque
$u''$ per aequationem $U''=0$, quae si supponitur esse gradus $m'''^{\mathrm{ti}}$, $m'''$ erit numerus
formae $2^{n-3}i'''$. Manifestum est, in serie aequationum $U=0$, $U'=0$,
$U''=0$ etc. $n^{\mathrm{tam}}$ fore gradus imparis adeoque radicem realem habere. Statuemus
brevitatis gratia $n=3$, ita ut aequatio $U''=0$ radicem realem $u''$ habeat,
nullo enim negotio perspicitur, pro quovis alio valore ipsius $n$ idem ratiocinium
valere. Tunc coëfficientem $M''$ per $u''$ et coëfficientes in $U'$ (quos fore functiones
integras coëfficientium in $Z$ facile intelligitur), sive per $u''$ et coëfficientes in
